% ========================================================
% SPEEDUP CON MEDIA RECORTADA
% ========================================================
\section{Speedup con media recortada}\label{sec:speedup}

\subsection{Protocolo de medición}

El protocolo está fijado en \texttt{tp/kmeans/cmd/benchmark} y no se decide en tiempo de corrida,
que es lo que evita que dos ejecuciones \enquote{según el mismo criterio} den cifras distintas:

\begin{itemize}
  \item \textbf{\cifra{repeticiones} repeticiones medidas por configuración}, más una ronda de
        calentamiento que se descarta.
  \item \textbf{Orden aleatorizado por bloques}: en cada ronda se barajan las configuraciones. No es
        un detalle: la caché de disco y el turbo del procesador premian a la configuración que corre
        segunda, así que medir siempre \enquote{primero secuencial, después concurrente} infla el
        speedup.
  \item \textbf{Media recortada simétrica al \cifra{recorte} por extremo}: se ordenan las
        \cifra{repeticiones} mediciones, se descartan las dos menores y las dos mayores, y se
        promedian las restantes. Descarta las corridas afectadas por un pico del sistema operativo
        sin elegir a mano cuáles.
  \item \textbf{El speedup es el cociente de las medias recortadas},
        $S = \bar{T}_{\text{sec}} / \bar{T}_{\text{conc}}$, y no el promedio de los cocientes por
        corrida, que sesga hacia arriba.
  \item \textbf{Intervalo de confianza del 95\,\% por bootstrap} de \cifra{bootstrap} réplicas
        sobre las dos muestras de tiempos: se remuestrean ambas con reemplazo, se recalcula el
        cociente de medias recortadas y se toman los percentiles 2,5 y 97,5.
  \item \textbf{Se mide solo el clustering.} La carga del CSV y la inicialización quedan fuera y se
        reportan aparte (Sección~\ref{sec:recursos}); ambas versiones parten del mismo archivo de
        centroides iniciales.
  \item \textbf{Dos experimentos independientes}: trabajo fijo (\cifra{iter-fijas} iteraciones,
        tolerancia cero), que compara el mismo trabajo; y hasta convergencia (tope de 60
        iteraciones, tolerancia $10^{-9}$), que compara el tiempo hasta la misma solución.
\end{itemize}

Las corridas oficiales se hicieron en el servidor \texttt{\cifra{maquina}} del equipo,
\cifra{cpus-gorgo} núcleos sin SMT, ocioso durante la medición; el benchmark completo tomó
\cifra{duracion-min}~minutos. Los tiempos crudos de cada repetición se guardan junto a los
agregados en \texttt{tp/reports/}, para poder recalcular cualquier estadística sin volver a medir.
Todas las tablas de esta sección y las siguientes se generan desde esos JSON con
\texttt{tp/scripts/tablas\_informe.py}; ninguna cifra se transcribió a mano.

\subsection{Tabla de speedup}\label{sec:fuerte}

% Generado por tp/scripts/tablas_informe.py desde tp/reports/*.json. No editar a mano.
\begin{table}[H]
\caption{Escalamiento fuerte con trabajo fijo: dataset completo}\label{tab:speedup-fuerte}
\footnotesize
\begin{tabular}{@{}rlrrlrr@{}}
\toprule
\textbf{$P$} & \textbf{Versión} & \textbf{Media rec. (ms)} & \textbf{Desv. (ms)} & \textbf{Speedup [IC 95\,\%]} & \textbf{Efic.} & \textbf{Karp--Flatt} \\
\midrule
--- & secuencial & 1.106,3 & 5,0 & --- & --- & --- \\
1 & concurrente & 1.141,2 & 6,1 & 0,97$\times$ [0,97; 0,97] & 97\,\% & --- \\
2 & concurrente & 574,4 & 3,4 & 1,93$\times$ [1,92; 1,93] & 96\,\% & 0,038 \\
4 & concurrente & 291,0 & 5,1 & 3,80$\times$ [3,77; 3,82] & 95\,\% & 0,017 \\
8 & concurrente & 190,1 & 6,2 & 5,82$\times$ [5,73; 5,89] & 73\,\% & 0,054 \\
16 & concurrente & 177,1 & 4,6 & 6,25$\times$ [6,17; 6,31] & 39\,\% & 0,104 \\
\bottomrule
\end{tabular}

{\footnotesize\textit{Nota.} Dataset completo, 2.831.486 viajes, trabajo fijo de 10 iteraciones. fischl-ubuntusv · go1.26.7 · 11 CPUs lógicas · 20 repeticiones medidas + 1 de calentamiento · media recortada al 10\,\% por extremo · IC del speedup por bootstrap de 2.000 réplicas · $k=8$ · chunk 16.384. Media rec. es la media recortada; la eficiencia es speedup/$P$. Karp--Flatt es la fracción serial efectiva $e = (1/S - 1/P)/(1 - 1/P)$; no está definida con $P = 1$.\par}
\end{table}


La Tabla~\ref{tab:speedup-fuerte} es el resultado principal. Sobre los \cifra{n}~viajes, la
versión secuencial tarda \cifra{seq-ms}~ms en \cifra{iter-fijas} iteraciones; la concurrente
alcanza \cifra{speedup-p2}$\times$ con dos workers, \cifra{speedup-p4}$\times$ con cuatro y
\cifra{speedup-p8}$\times$ con ocho. Los intervalos son estrechos, del orden de centésimas, porque
la desviación entre repeticiones es de pocos milisegundos sobre cientos: el protocolo produce
mediciones estables.

Dos lecturas que conviene no confundir. Con $P = 1$ la versión concurrente es un
\cifra{sobrecosto-p1} más lenta que la secuencial, y el intervalo no toca el 1: es el precio fijo
del pool, del troceado y de la reducción, que el paralelismo tiene que recuperar antes de ganar algo.
Y con $P = 16$ sobre \cifra{cpus-gorgo} núcleos el speedup todavía sube un poco
(\cifra{speedup-p16}$\times$), porque hay más bloques listos para ocupar cualquier núcleo que se
libere, pero la eficiencia por worker cae a \cifra{eficiencia-p16}: no son 16 procesadores, es
sobresuscripción.

\subsection{El mismo resultado hasta convergencia}

% Generado por tp/scripts/tablas_informe.py desde tp/reports/*.json. No editar a mano.
\begin{table}[H]
\caption{Escalamiento fuerte hasta convergencia: dataset completo}\label{tab:speedup-convergencia}
\footnotesize
\begin{tabular}{@{}rlrrlrr@{}}
\toprule
\textbf{$P$} & \textbf{Versión} & \textbf{Media rec. (ms)} & \textbf{Desv. (ms)} & \textbf{Speedup [IC 95\,\%]} & \textbf{Efic.} & \textbf{Karp--Flatt} \\
\midrule
--- & secuencial & 6.016,1 & 14,0 & --- & --- & --- \\
1 & concurrente & 6.174,7 & 13,0 & 0,97$\times$ [0,97; 0,98] & 97\,\% & --- \\
2 & concurrente & 3.107,8 & 8,0 & 1,94$\times$ [1,93; 1,94] & 97\,\% & 0,033 \\
4 & concurrente & 1.579,8 & 9,6 & 3,81$\times$ [3,80; 3,82] & 95\,\% & 0,017 \\
8 & concurrente & 1.043,1 & 7,4 & 5,77$\times$ [5,75; 5,79] & 72\,\% & 0,055 \\
16 & concurrente & 982,0 & 16,2 & 6,13$\times$ [6,08; 6,17] & 38\,\% & 0,107 \\
\bottomrule
\end{tabular}

{\footnotesize\textit{Nota.} Dataset completo, 2.831.486 viajes, hasta convergencia; ambas versiones pararon en 60 iteraciones (tope del protocolo). fischl-ubuntusv · go1.26.7 · 11 CPUs lógicas · 20 repeticiones medidas + 1 de calentamiento · media recortada al 10\,\% por extremo · IC del speedup por bootstrap de 2.000 réplicas · $k=8$ · chunk 16.384. Media rec. es la media recortada; la eficiencia es speedup/$P$. Karp--Flatt es la fracción serial efectiva $e = (1/S - 1/P)/(1 - 1/P)$; no está definida con $P = 1$.\par}
\end{table}


El segundo experimento (Tabla~\ref{tab:speedup-convergencia}) deja correr ambas versiones hasta el
criterio de parada. Con la tolerancia de $10^{-9}$ ninguna converge antes del tope de 60 iteraciones,
así que en la práctica compara seis veces más trabajo que el primero, y da lo mismo dentro del ruido:
\cifra{speedup-conv-p8}$\times$ contra \cifra{speedup-p8}$\times$ con ocho workers. Eso descarta que
el resultado dependa de cuál de los dos experimentos se elija.

\subsection{Speedup según el tamaño del problema}

% Generado por tp/scripts/tablas_informe.py desde tp/reports/*.json. No editar a mano.
\begin{table}[H]
\caption{Speedup según el tamaño del problema y la cantidad de workers}\label{tab:speedup-tamano}
\small
\begin{tabular}{@{}rrrrrrr@{}}
\toprule
\textbf{$n$ (viajes)} & \textbf{Secuencial (ms)} & \textbf{$P=1$} & \textbf{$P=2$} & \textbf{$P=4$} & \textbf{$P=8$} & \textbf{$P=16$} \\
\midrule
1.000 & 0,5 & 0,89$\times$ & 0,90$\times$ & 0,94$\times$ & 0,92$\times$ & 0,86$\times$ \\
10.000 & 4,4 & 0,96$\times$ & 0,97$\times$ & 0,93$\times$ & 0,95$\times$ & 0,95$\times$ \\
100.000 & 41,8 & 0,97$\times$ & 1,88$\times$ & 2,87$\times$ & 3,69$\times$ & 3,64$\times$ \\
1.000.000 & 388,0 & 0,97$\times$ & 1,89$\times$ & 3,62$\times$ & 5,49$\times$ & 5,93$\times$ \\
2.831.486 & 1.106,3 & 0,97$\times$ & 1,93$\times$ & 3,80$\times$ & 5,82$\times$ & 6,25$\times$ \\
\bottomrule
\end{tabular}

{\footnotesize\textit{Nota.} Trabajo fijo de 10 iteraciones. Cada celda es el speedup puntual de la media recortada; los intervalos de las filas del dataset completo están en la Tabla~\ref{tab:speedup-fuerte}.\par}
\end{table}


La Tabla~\ref{tab:speedup-tamano} repite el experimento de trabajo fijo sobre submuestras
sistemáticas del gold. El speedup crece con $n$: con mil o diez mil viajes la versión concurrente
pierde con cualquier $P$, con cien mil ya gana pero la eficiencia es baja, y a partir de un millón
la curva se parece a la del dataset completo. Esa dependencia del tamaño es la que la
Sección~\ref{sec:equilibrio} mide con más resolución para ubicar el punto de equilibrio.
